\section{Interacting with the ledger in Axi}

In this appendix, we will see how Axi programs can interact with Khalani's ledger. We will start with an overview of the ``life cycle'' of resources, which is a mental model of what can happen to resources (and how) both when they are on the ledger and when they are being used by a program. Then we will look at the primitive (built-in) types and functions that Axi uses to interact with the ledger. Finally, we will discuss a slightly longer example of an on-chain service, together with example transactions that set it up and use it.

\subsection{The resource life cycle}

Resource types can be divided into built-in and user-defined. The built-ins include gas, the native coin used to pay transaction fees. The built-ins exist right from the genesis block. A new user-defined resource type is created by a transaction that deploys a package of Axi code that contains its definition. Because packages live on-chain forever and cannot be deleted, a resource type, once it comes into existence, will also exist forever.

The life cycle of a particular resource, i.e. an element of some resource type, consists of three stages:

\begin{itemize}
  \item \textbf{Creation}. At some point, the resource is created, in accordance with the interface of the package that contains the definition of the resource's type. After creation, the resource is placed on the ledger either in an owned cell or in a shared cell. Note that it can be placed in the cell either on its own or as a part of a larger, compound resource.
  \item \textbf{Transactions}. The resource participates in transactions. Resources in owned cells can be freely transferred by their owner. Resources in shared cells can be interacted with via the interface provided by the defining package of the cell's type.
  \item \textbf{Destruction (optional)}. If the defining package permits it, the resource might at some point be destroyed. Examples of destructible resources are stablecoins that are backed by an off-chain reserve, where a coin is destroyed when its off-chain collateral is removed from the treasury, and a wrapper for an asset bridged from another blockchain, which is destroyed when the external asset ceases to be bridged.
\end{itemize}

The most important part of the life cycle, in which most resources will spend almost all their lives, is the phase when they can participate in transactions. An (ordinary) transaction consists of an arbitrary Axi program that transforms the state of the ledger. This program, when executed, interacts with the ledger in three steps:

\begin{itemize}
  \item \textbf{Withdraw}. First, resources are ``withdrawn'' from the ledger into the program. The ledger's ownership model ensures the safety of this step, for example by making sure that the program cannot withdraw resources that don't belong to the transactor.
  \item \textbf{Transform}. Second, the program transforms its resources by performing various computations, exchanging the resources for other resources and moving them around, etc. The safety of this step is guaranteed by the type system. The resources cannot be unlawfully destroyed or copied because the type system tracks usage. Nonsensical operations cannot be performed on them, because the program is well-typed. They cannot get lost in complex pattern matching, because exhaustiveness is checked, and so on.
  \item \textbf{Deposit}. Third, the program ``deposits'' the transformed resources back to the ledger. The safety of this step is partially ensured by the type system, which checks that all resources held by the program are deposited back. However, since resources can be freely transferred, it's up to the transactor to make sure that they are transferred to the correct address and not to somebody else.
\end{itemize}

As we can see, the resource life cycle is determined by two orthogonal but complementary forces: package and service interfaces and the ownership system. The creator of the module that defines a new resource type sets the interface for all uses of all resources of this type. For a service (i.e. a shared mutable cell), its creator defines, through its interface, how the service handles all its resources. For an owned cell, its owner determines how the ownership of the resource changes, with which services it will interact (and through which methods of these services), and what primitive operations (coming from the defining package) will be performed on the resource.

\subsection{Ledger primitives}

In this section, we will see what Khalani's ledger looks like from Axi's perspective. To do this, we will discuss the primitive (i.e. built-in) types and functions that Axi programs use to interact with the ledger.

\lstinputlisting{listings/Primitives/Ledger.axi}

\code{Ledger} is the type that represents (the current state of) the ledger. When we said that the body of an ordinary transaction is an ``arbitrary Axi program that transforms the ledger'', we meant precisely that it is a function of type \code{Ledger -> Ledger}.

Every function that interacts with the ledger takes (the current version of) the ledger as input and also returns (the new version of) the ledger as output. In general, the type of such a function is \code{A -> Ledger -> B * Ledger}, i.e. it takes a value of some type \code{A} and the ledger, and returns a value of some type \code{B} and the (updated) ledger.

The transaction body will, in general, be a composition of many such functions. It will get the initial input of type \code{Ledger} from the transaction context and then pass it to the first of these functions, which will in turn pass it to the next one, and so on. Because \code{Ledger} is a linear type, the type system guarantees that values of this type must be used exactly once and cannot be copied or deleted. This means that the ledger is acted upon and then passed along until the transaction body returns the final modified version of it. This is how all interactions with the ledger are represented in code.

\lstinputlisting{listings/Primitives/transactor.axi}

An example function that needs the ledger is \code{transactor}, which consumes the old ledger and returns a reference to the account which originated the current transaction (i.e. the transactor) together with a new ledger.

\lstinputlisting{listings/Primitives/Addresses.axi}

\code{Address} is the data type of addresses, including both accounts and cell addresses. \code{Account} is the type of addresses that are known to be addresses of accounts. \code{Ref} is a type constructor that turns resource types into data types that represent ``references'' to cells that contain resources of this type. The problem we want to solve is that, when interacting with shared cells, we don't know whether a shared cell with the given address exists, so we would need to check that first and handle the possible errors, but doing that over and over is tedious. As a solution, we perform a kind of ``handshake'' with the cell using the function \code{referenceof} which, if a cell with the given address exists, returns a reference to it. In practice, \code{Ref A} is just a wrapper around an \code{Address} that points to a shared cell that holds a value of type \code{A}. Besides shared cells, we also use \code{Ref} with \code{Account}s, which, due to account abstraction, also have their own cells. This helps make the type signature of \code{new} simpler.

The function \code{addressof} returns the address when given the reference. The function \code{referenceof}, when given an address, together with the ledger, returns an optional reference, together with the ledger. If such a cell doesn't exist, the result is \code{none} (together with the new ledger). If it does, the result is \code{some ref} (together with the new ledger), where \code{ref} is the reference to the cell.

\lstinputlisting{listings/Primitives/new.axi}

\code{new} is a function that takes an account reference (the type argument is implicit), together with the ledger, and creates a new cell on the ledger. It returns the address of the created cell together with a function of type \code{A -> Ledger}. This function is a ``setter'' for the cell -- when it is given a value of type \code{A}, it sets the cell contents to this value and then returns the new ledger. Note that the type \code{A -> Ledger} is linear, so the type system makes sure that it is called exactly once. Because of this, the cell can't be left empty, nor can it be overwritten right after being created. This is a pattern that we will also see when creating and updating the shared cells.

\lstinputlisting{listings/Primitives/delete.axi}

\code{delete} is a function that takes a type \code{A} (note that the symbol \code{@} marks that this type argument must be given explicitly) and a cell address, together with the ledger, and returns an optional \code{A}, together with the new ledger. If the transactor is not the owner of the cell, or the value stored in the cell is not of type \code{A}, then the operation fails and the result is \code{none} (together with the new ledger). Otherwise, the operation succeeds, the cell is destroyed and the result is \code{some a} (together with the new ledger), where \code{a} is the value that was stored in this cell.

\lstinputlisting{listings/Primitives/share.axi}

\code{share} is a function that takes just the ledger (the type argument is implicit), creates a new mutable shared cell, and returns a reference to the created cell, together with a setter function that, when given a value of type \code{A}, sets the cell's contents to this value and returns the new ledger.

\lstinputlisting{listings/Primitives/update.axi}

\code{update} is a function that takes a reference to a mutable shared cell, together with the ledger, and returns a value of some type \code{A}, together with a setter function that, when given a value of type \code{A}, sets the cell's contents to this value and returns the new ledger. Because the type \code{A} is linear, we must use the returned value exactly once, and because the type \code{A -> Ledger} is also linear, we must call the setter with some \code{A}. In practice, this means that if the type \code{A} is abstract, all we can do is modify the value of type \code{A} that we were given and then put it back into the cell by calling the setter.

\lstinputlisting{listings/Primitives/freeze.axi}

\code{freeze} is a function that takes just the ledger (the data type argument \code{A} is implicit), creates a new immutable shared cell, and returns the address of the created cell, together with a setter function that, when given a value of type \code{A}, sets the cell's contents to this value and returns the new ledger.

\lstinputlisting{listings/Primitives/peek.axi}

\code{peek} is a function that takes a type \code{A} and an address, together with the ledger, reads the contents of an immutable shared cell at this address, and returns an optional \code{A}, together with the new ledger. If the address does not point to an immutable shared cell, or the value stored in that cell is not of type \code{A}, then the operation fails and the result is \code{none} (together with the new ledger). Otherwise the operation succeeds and the result is \code{some a} (together with the new ledger), where \code{a} is the data stored in the cell.

\subsection{A longer example -- donut shop}

In this section, we will see a longer example that demonstrates a simple on-chain service -- the donut shop. This example was inspired by a similar example from the \href{https://docs.sui.io/concepts/object-ownership/shared}{Sui documentation}.

\begin{example}{Donut Shop}

\lstinputlisting{listings/DonutShop.axi}

\end{example}

The above listing shows the code of the module \code{DonutShop}. This module provides functions for creating, managing and interacting with donut shops, which will live on-chain as \textit{services}. We use the name service to refer to a mutable shared cell that stores a value whose type is abstract, i.e. whose definition is not exported by the module that defines it. Because of this, the only way to interact with a service is to get it from the ledger, call some of its methods (i.e. functions from the module that defines it) and then put it back into its cell.

When the code above gets deployed onto the chain, the resulting package will effectively become a public shared library that everybody can use in their own code. However, to make the examples simpler, we suppress details related to packages and instead pretend that it's the module itself that gets deployed onto the chain.

First, on lines 7-11, we indicate which functions we want the module to export. These functions become the interface of the module through which other code interacts with types from this module. We export all the functions we define, but we export neither the types nor their constructors or destructors, because we want to keep all of them abstract, so that elements of the \code{DonutShop} type can become on-chain services.

On lines 14-15 we define the \code{Donut} type. Because of the annotation \code{:\ Type?}, \code{Donut} is an affine type, i.e. its elements must be used at most once. Every donut has a serial number -- this guarantees that every shop has an infinite supply of donuts, but otherwise this detail isn't very relevant to the example.

On lines 20-24 we define the \code{DonutShop} type. A shop stores its own on-chain address (so that it can check if the person trying to collect its profits is its owner), its balance (the coins it earned so far from selling donuts), the price of a single donut and the serial number of the next donut it will sell. \code{DonutShop} is a linear type, so every shop must be used exactly once. This means that functions which interact with the donut shop will take the shop as input and return an updated shop as output.

On lines 27-28 we define the \code{DonutShopOwnershipToken} type. Such a token contains a reference to the on-chain cell in which the shop lives and acts as a ``key'' that allows the shop owner (and only him) to withdraw the profits made by the shop. We expect such ownership tokens to become a common programming pattern in Khalani. We will explain how these tokens work later.

On lines 31-43 we define the function \code{init}. It takes a donut price, together with the ledger, sets up a new donut shop on the ledger, and returns its ownership token, together with the new ledger. It works as follows. First, on line 34, we create a new mutable shared cell that will hold the shop. In return we get a reference and a setter. Then on lines 35-36 we define a new ownership token using the reference. Then on lines 37-41 we define a new shop using the reference's address. Finally, on line 43, we return the ownership token, together with the new ledger obtained by calling the setter \code{publishShopToLedger} on the new shop.

On lines 47-62 we define the function \code{buyDonut}. It takes a coin used as payment and a donut shop, and returns an optional donut (depending on whether the purchase succeeds or not), the change left after paying, and the updated donut shop. It works as follows. First, on line 49, we split the payment coin to see if there's enough of it to pay the donut's price. On line 50 we check if the split succeeded. If not (line 51), we return \code{none}, the change (which is equal to the original payment coin) and the unmodified shop. If the split succeeded (line 52), we create a new donut (line 55), compute the shop's new balance (line 56) and update the shop (lines 57-60) to reflect the new balance and new serial number. Finally, on line 62, we return the donut, the change and the updated shop.

Note that, even though on line 49 we access the \code{shop}'s \code{price} field, this does not count as a usage of \code{shop} -- this is because the accessed field, \code{price}, is data and not a resource. We allow accessing data fields of records without consuming the record in order to avoid lengthy decomposition and recomposition of the record every time we need to peek at some data. On line 56, when we access the \code{shop}'s \code{balance} field, which is a resource, it does count as a usage of \code{shop}. However, when on line 58 we use \code{shop} as part of the \textit{record update syntax} to build the updated shop, this again does not count as a usage, because all the fields that still ``remain'' in \code{shop} are data fields. \code{newShop} has the same fields as \code{shop}, except for the fields that were overwritten on lines 59-60. Because on line 59 we provide the \code{balance} field, which was ``withdrawn'' from \code{shop} by the field access on line 56, the definition of \code{newShop} through record update of \code{shop} is valid and follows correct resource usage.

On lines 65-74 we define the function \code{collectProfits}. It takes an ownership token and a shop and returns coins, the new shop and the new token. It works as follows. On line 67 we check if the address of the ownership token's reference is the same as the address of the shop. If yes, on lines 69-70 we withdraw the profits from the shop and update the shop's balance to zero, and then on line 72 we return the profits, the ownership token and the updated shop. If the token's address is different from the shop's address, the operation is unauthorized, and in such a case on line 74 we return a zero-value coin, the ownership token and the original shop.

Last but not least, on lines 77-79 we define the function \code{getMyShopRef} which, when given an ownership token, returns a reference to the shop and the ownership token. This function is needed because, even though inside the \code{DonutShop} module the ownership token's reference can be accessed using record field access syntax, outside the module this is not the case (because the definition of \code{DonutShopOwnershipToken} is not exported). Thus, without \code{getMyShopRef}, the shop owner couldn't get a reference to his shop from his ownership token. This wouldn't be a fatal flaw, but only a mild annoyance.

These technical descriptions of the functions' inner workings might sound scary, but in practice things are quite simple. Anybody can call \code{init} to set up a new donut shop. In return he gets an ownership token, which allows him to collect the shop's profits and also contains the shop's reference (and thus address), so that he remembers where his shop is located. Anybody can try to buy a donut from any donut shop by calling the \code{buyDonut} function. If the argument contains enough coins, he will get his donut and change. If not, he doesn't get a donut but gets his coins back. In both cases, he can continue his interaction with the shop. Finally, anybody with any donut shop ownership token can try to collect profits from any donut shop. If the token happens to be the correct token for the particular shop, he will get his profits. If not, he gets a zero-value coin. In both cases, he gets his ownership token back and can continue his interaction with the shop.

Now, why does the donut shop ownership token work? Why can't anybody steal the shop's profits? First, the \code{DonutShopOwnershipToken} is abstract, i.e. its definition is not exported by the module, so the only way to create such a token is by calling \code{init}. Second, this type is linear, so that the ownership token cannot be copied or deleted (although it might be sent to a dummy address, becoming effectively lost forever). This ensures that on-chain there always exists a unique ownership token for every donut shop. Because every donut shop stores its own on-chain address, and this address cannot be changed either by moving the shop on-chain or by modifying the shop's \code{thisShop} field, the check that the \code{collectProfits} function performs is enough to identify whether the caller is the shop's rightful owner or not. This argument could in principle be turned into a theorem, but we leave this for further work.

Another interesting property of this design is that, although ownership tokens cannot be forged, copied or deleted, they can be exchanged, bought and sold. This means, for example, that a donut shop owner could sell his shop for a lump sum, letting the buyer collect the profits over time. Thus, we can see the donut shop ownership token as a very simple kind of stock that can be traded OTC and pays dividends on demand, but in a single lump sum (the owner can't withdraw only a part of the profits). These stock-like properties were not explicitly designed by us when we were coming up with this example. Instead, they arose naturally from the ``ownership token'' programming pattern, abstract types and linear typing.

\begin{example}{A transaction that sets up a new donut shop}

\lstinputlisting{listings/DonutShopInit.axi}

\end{example}

The above example shows a transaction that sets up a new donut shop. On lines 1-2, we import the \code{init} function from the \code{DonutShop} module. On lines 4-14 we define the \code{main} function, which describes how the transaction will modify the ledger once it is executed. On line 8 we ask for a reference to our own account (so that we can create new cells later). On line 10 we call \code{init} to set up a new donut shop in which each donut costs 1000 \code{KhalaniCoin}s. On line 11 we create a new cell owned by us that will hold the shop's ownership token. On line 12 we call the setter to put our token into this cell. Finally, on line 14, we return the ledger.

\begin{example}{A transaction that buys a donut}

\lstinputlisting{listings/DonutShopBuy.axi}

\end{example}

The above example shows a transaction that buys a donut. On lines 1-3 we import the types \code{Donut}, \code{DonutShop} and the function \code{buyDonut} from the \code{DonutShop} module. On lines 5-7, we define a constant that holds the address of our cell with coins that we will use to pay for the donut. On lines 9-11 we define a constant that holds the address of the donut shop from which we will buy the donut. On lines 13-41 we define the \code{main} function. On line 17 we ask for a reference to our own account. On line 18 we ask for a reference to the donut shop. On line 20 we check if there is indeed a donut shop at the given address. If not (line 21), then we finish by returning the ledger. If yes, we now have a reference to the shop (line 22). On line 23 we attempt to withdraw our coins from their cell and on line 24 we check if it succeeded. If not (for example because we made a typo in the address), we finish by returning the ledger (line 25). If yes, we now have the coins (line 26). On line 28 we retrieve the shop. On line 29 we call \code{buyDonut} on the shop, trying to pay with our coins. On line 30 we put the shop back into its cell. On line 31 we create a new cell to hold the change that was left after paying for the donut. On line 32 we put the change into this cell. On line 34 we check if we bought the donut successfully. If not (for example because we didn't have enough money), then we finish by returning the ledger (line 35). If yes, we now have the donut (line 36). On line 38 we create a new cell owned by us that will hold the donut. On line 39 we put the donut into this cell. On line 41 we return the final ledger.

\begin{example}{A transaction that tries to collect profits from a donut shop}

\lstinputlisting{listings/DonutShopCollect.axi}

\end{example}

The above example shows a transaction that tries to collect profits from a donut shop. On lines 1-4 we import the necessary types and functions from the module \code{DonutShop}. On lines 6-8 we define a constant that holds the address of our donut shop ownership token. On lines 10-31 we define the \code{main} function. On line 14 we ask for a reference to our own account. On lines 15-16 we try to pick up our donut shop ownership token from the ledger. On line 18 we check if we got the token. If not (for example, because this cell holds something else and not our token), we finish by returning the ledger (line 19). If yes, we now have the ownership token (line 20). On line 22 we extract the shop reference from the token. On line 23 we pick up the shop from the ledger. On line 24 we call \code{collectProfits} on the shop. On line 25 we put the shop back into its cell. On line 26 we create a new cell owned by us that will hold the profits. On line 27 we put the profits into this cell. On line 28 we create a new cell owned by us that will hold our token (recall that the old cell was destroyed when we picked up the token). On line 29 we put the ownership token into this cell. On line 31 we return the final ledger.
